% Shared result, cited from solutions with \appendixref{field-norm}.
% Moved on 2026-09-30 from the appendix of problem 0041 (A.2 there), text unchanged: only the links to other
% results became \appendixref.  Earlier version: spring 2025 problem set 9, solutions appendix (static/uploads/problems/spring_2025/problemset_9, commit 7bdae2b^), A.2 there.
% The equation (**) below is used by problem 0041: its label has the global name appendix:field-norm:A.2.3.
\begin{appendixitem}{Field norm}
For a field extension \( K \hookrightarrow L \), for each \( \alpha \in L \), we can consider the \( K \)-linear endomorphism of mutliplication by \(\alpha\):
\[
[\times \alpha] : L \rightarrow L.
\]
If the extension is \(K\)-finite-dimensional (then it is algebraic), we can compute its characteristic polynomial \( P_{\operatorname{char}, [\times \alpha]}(X) \in K[X] \). In particular, we can compute its determinant (the constant term of \( P_{\operatorname{char}, [\times \alpha]}(X) \)), which defines the \textit{field norm}:
\[
N_{L/K} : L \rightarrow K, \quad N_{L/K}(\alpha) := {\det}_{K}([\times \alpha]).
\]
Since \( \forall \beta \in L,\ [\times \alpha \beta] = [\times \alpha] \circ [\times \beta] \), we have
\[
N_{L/K}(\alpha \beta) = N_{L/K}(\alpha) N_{L/K}(\beta).
\]
In the case where \( L/K \) is Galois (which is equivalent to \( K \hookrightarrow L \) being a normal and separable field extension), we must have that the roots in \( L \) of the minimal polynomial of \( \alpha \) in \( K[X] \) satisfy
\[
\left\{ \sigma(\alpha) \mid \sigma \in \operatorname{Gal}(L/K) \right\} = \operatorname{Root}_{P_{\min, \alpha}(X)}(L) = \operatorname{Root}_{P_{\min, \alpha}(X)}\left( K^{\mathrm{alg}} \right),
\]
where the first equality follows from standard Galois theory, and the second from the fact that the extension is normal. Now, because \( P_{\min, \alpha}(X) = P_{\min, [\times \alpha]}(X) \) (a mental exercise), we have:
\[
\operatorname{Root}_{P_{\operatorname{char}, [\times \alpha]}(X)}\left( K^{\mathrm{alg}} \right) = \operatorname{Root}_{P_{\min, [\times \alpha]}(X)}\left( K^{\mathrm{alg}} \right) = \operatorname{Root}_{P_{\min, \alpha}(X)}\left( K^{\mathrm{alg}} \right),
\]
where the first equality follows since the roots (in \( K^{\mathrm{alg}} \)) of the minimal polynomial of an endomorphism (on a finite-dimensional vector space) are precisely its eigenvalues. Combining:
\begin{equation}
\label{A.2.1}
\tag{1}
\operatorname{Root}_{P_{\operatorname{char}, [\times \alpha]}(X)}\left( K^{\mathrm{alg}} \right) = \left\{ \sigma(\alpha) \mid \sigma \in \operatorname{Gal}(L/K) \right\}.
\end{equation}
If \( \alpha \) generates the extension, i.e. \( L = K(\alpha) \), or equivalently \( \deg\left( P_{\min, \alpha}(X) \right) = [L:K] \), then, since the extension is finite and Galois, we have \( \left| \operatorname{Gal}(L/K) \right| = [L:K] \). From this, as \( P_{\min, \alpha}(X) \) is separable, it follows that for each \( \sigma, \sigma' \in \operatorname{Gal}(L/K) \), \( \sigma \neq \sigma' \) implies \( \sigma(\alpha) \neq \sigma'(\alpha) \). Thus:
\begin{equation}
\label{A.2.2}
\tag{2}
\left|\left\{ \sigma(\alpha) \mid \sigma \in \operatorname{Gal}(L/K) \right\}\right|=[L:K]=\deg\left(P_{\operatorname{char}, [\times \alpha]}(X)\right).
\end{equation}
From \eqref{A.2.1} and \eqref{A.2.2}, we obtain:
\[
P_{\operatorname{char}, [\times \alpha]}(X) = \prod_{\sigma \in \operatorname{Gal}(L/K)} \left( X - \sigma(\alpha) \right)\footnote{In fact, this holds even if \( \alpha \) does not generate the extension (but still \( L/K \) is finite and Galois), as one can show that the multiplicities of each root of the characteristic polynomial \( \theta \in \operatorname{Root}_{P_{\operatorname{char}, [\times \alpha]}(X)}\left( K^{\mathrm{alg}} \right) \) coincide with the number of automorphisms of \( L \) fixing \( K \) and sending \( \alpha \) to \( \theta \), that is, \( \left| \left\{ \sigma \in \operatorname{Gal}(L/K) \mid \sigma(\alpha) = \theta \right\} \right| \).},
\]
and so, by Vieta's formula:
\begin{equation}
\label{appendix:field-norm:A.2.3}
\tag{**}
N_{L/K}(\alpha) = \prod_{\sigma \in \operatorname{Gal}(L/K)} \sigma(\alpha).
\end{equation}
\end{appendixitem}
